
\subsection*{Preserving Versus Reflecting Order}

\begin{definitionbox}{Definition (Order-Reflecting Map)}
\begin{definition}[Order-Reflecting Map]\label{def:order-reflecting-map}
Let $(A,\leq)$ and $(B,\leq')$ be ordered sets. A function $f:A\to B$ is
\emph{order-reflecting} if
\[
f(x)\leq' f(y)\quad\Longrightarrow\quad x\leq y
\]
for all $x,y\in A$.
\end{definition}
\end{definitionbox}

\begin{remark*}[Standard quantified statement]
\[
  \forall x,y\in A,\quad
  f(x)\leq' f(y)\to x\leq y.
\]
\end{remark*}

\begin{remark*}[Predicate reading]
\[
  \operatorname{OrderReflecting}(f,\leq,\leq')
  \iff
  \forall x,y\in A,\ (f(x)\leq' f(y)\to x\leq y).
\]
\end{remark*}

\begin{remark*}[Interpretation]
An order-reflecting map pulls output comparisons back to input comparisons.
\end{remark*}

\begin{remark*}[Exposition]
Order-preserving pushes an inequality forward. Order-reflecting pulls an
inequality back. Together they say
\[
x\leq y \quad\Longleftrightarrow\quad f(x)\leq' f(y).
\]
\end{remark*}

\begin{dependencies}
\begin{itemize}
  \item \hyperref[def:function]{Function}
  \item \hyperref[def:partially-ordered-set]{Partially Ordered Set}
\end{itemize}
\end{dependencies}

\begin{definitionbox}{Definition (Order Embedding, Function-Property Form)}
\begin{definition}[Order Embedding, Function-Property Form]\label{def:order-embedding-function-form}
An \emph{order embedding} is a function that is both order-preserving and
order-reflecting.
\end{definition}
\end{definitionbox}

\begin{remark*}[Standard quantified statement]
\[
  \operatorname{OrderEmbedding}(f)
  \iff
  \operatorname{OrderPreserving}(f)\wedge\operatorname{OrderReflecting}(f).
\]
\end{remark*}

\begin{remark*}[Interpretation]
An order embedding preserves and reflects order, so it identifies the domain
with an ordered copy of its image.
\end{remark*}

\begin{remark*}[Exposition]
This is the same notion as \hyperref[def:order-embedding]{Order Embedding} in
the orderings chapter, restated here to emphasize the function-property
decomposition.
\end{remark*}

\begin{dependencies}
\begin{itemize}
  \item \hyperref[def:order-preserving-map-ordered-sets]{Order-Preserving Map}
  \item \hyperref[def:order-reflecting-map]{Order-Reflecting Map}
\end{itemize}
\end{dependencies}

\begin{propositionbox}{Proposition (Order Embeddings Are Injective, Revisited)}
\begin{proposition}[Order Embeddings Are Injective, Revisited]\label{prop:order-embedding-injective-revisited}
Every order embedding between partial orders is injective.
\hyperref[prf:order-embedding-injective-revisited]{Go to proof.}
\end{proposition}
\end{propositionbox}

\begin{remark*}[Standard quantified statement]
\[
  \operatorname{OrderEmbedding}(f)\to \operatorname{Injective}(f).
\]
\end{remark*}

\begin{remark*}[Interpretation]
Order embeddings between partial orders cannot collapse two distinct inputs.
\end{remark*}

\begin{remark*}[Exposition]
An order-preserving map can collapse distinct comparable elements. An
order-reflecting map cannot: if $f(x)=f(y)$, then both $f(x)\leq' f(y)$ and
$f(y)\leq' f(x)$ hold. Reflection gives $x\leq y$ and $y\leq x$, and
antisymmetry gives $x=y$.
\end{remark*}

\begin{dependencies}
\begin{itemize}
  \item \hyperref[def:order-embedding-function-form]{Order Embedding, Function-Property Form}
  \item \hyperref[def:injective]{Injective Function}
  \item \hyperref[def:antisymmetric]{Antisymmetric Relation}
\end{itemize}
\end{dependencies}

\subsection*{Isomorphisms Onto Images}

\begin{theorembox}{Theorem (Order Embedding Is Isomorphism Onto Image)}
\begin{theorem}[Order Embedding Is Isomorphism Onto Image]\label{thm:order-embedding-isomorphism-onto-image}
Let $f:(A,\leq)\to(B,\leq')$ be an order embedding. Then the corestriction
\[
f_{\operatorname{im}}:(A,\leq)\to(f(A),\leq'_{f(A)})
\]
is an order isomorphism.
\hyperref[prf:order-embedding-isomorphism-onto-image]{Go to proof.}
\end{theorem}
\end{theorembox}

\begin{remark*}[Standard quantified statement]
\[
  \operatorname{OrderEmbedding}(f)
  \to
  \operatorname{OrderIsomorphism}(f_{\operatorname{im}}).
\]
\end{remark*}

\begin{remark*}[Interpretation]
An order embedding becomes an order isomorphism after the codomain is restricted
to the image.
\end{remark*}

\begin{dependencies}
\begin{itemize}
  \item \hyperref[def:corestriction-to-image]{Corestriction to the Image}
  \item \hyperref[def:order-embedding-function-form]{Order Embedding, Function-Property Form}
  \item \hyperref[def:order-isomorphism-ordered-sets]{Order Isomorphism}
\end{itemize}
\end{dependencies}

\begin{theorembox}{Theorem (Inverse of an Order Isomorphism)}
\begin{theorem}[Inverse of an Order Isomorphism]\label{thm:inverse-order-isomorphism}
If $f:(A,\leq)\to(B,\leq')$ is an order isomorphism, then
$f^{-1}:(B,\leq')\to(A,\leq)$ is an order isomorphism.
\hyperref[prf:inverse-order-isomorphism]{Go to proof.}
\end{theorem}
\end{theorembox}

\begin{remark*}[Standard quantified statement]
\[
  \operatorname{OrderIsomorphism}(f)\to \operatorname{OrderIsomorphism}(f^{-1}).
\]
\end{remark*}

\begin{remark*}[Interpretation]
The inverse of an order isomorphism preserves the same order structure in the
reverse direction.
\end{remark*}

\begin{remark*}[Exposition]
A bijective order-preserving map should not be treated as an order isomorphism
until the inverse is known to be order-preserving, or equivalently until the
map is known to reflect order.
\end{remark*}

\begin{dependencies}
\begin{itemize}
  \item \hyperref[def:order-isomorphism-ordered-sets]{Order Isomorphism}
  \item \hyperref[def:inverse]{Inverse Function}
  \item \hyperref[def:order-reflecting-map]{Order-Reflecting Map}
\end{itemize}
\end{dependencies}
